% SGA 3，第 VII 篇_B，英文重建本——第 1 轮（中文译本）。
% 范围：依据原文件第 1--3 页，从标题至 0.1.3；止于原文件第 3 页的 0.2 之前。
% 合并阅读版页码： 514--516. 再版印刷页码： 503--505.
% 本范围内的原版页边定位： 477--480.
% 依据： Polo--Gille Exp7B-13oct24.pdf,
% SHA-256 D303690E0195654BBE7C3DAA5D89EA23E24CD709A3575217234AC19AE8836694.
% 曾参阅 Jacob Reinhold 的 CC BY 4.0 Markdown，作为对比版本来源。

\chapter*{第 \texorpdfstring{VII\textsubscript{B}}{VII-B} 篇\\
群概形的无穷小研究}
\phantomsection
\label{sga3:VIIB:title}
\addcontentsline{toc}{chapter}{第 VII-B 篇. 群概形的无穷小研究}
\markboth{第 VII\textsubscript{B} 篇．群概形的无穷小研究}%
  {第 VII\textsubscript{B} 篇．群概形的无穷小研究}

\begin{center}
  \textit{P. Gabriel 著}
\end{center}

\setcounter{footnote}{-1}
\section*{B) 形式群}
\phantomsection
\label{sga3:VIIB:B}
\addcontentsline{toc}{section}{B) 形式群}

% 依据： Exp7B 原文件页码 1; 合并阅读版页码 514; 印刷页码 503.
\noindent\footnote{编者注：2024 年 10 月 13 日版本。}
形式群的研究通常极为简单。如果读者从下面这些页中不能清楚地看出这一点，
那应当归咎于一位声称了解“域以外的某些东西”上的形式群的算术学家。\footnote{例如，Lubin 和 Tate 的工作（参见\ \textup{[LT65]}）表明了完备
Noether 局部环上的形式群的重要性。关于任意基上的形式群，以及提升与形变问题的研究，
特别是关于 Barsotti--Tate 群（“\(p\)-可除群”）的研究，已经产生了大量文献；
例如参见 \textup{[LT66, Ta67, Gr74, Me72, La75, Fo77, Il85, Br00]}。
特别地，本篇的结果大多在 \textup{[Fo77]} 第~I 章中重新论述。}
因此，我们对“在 Artin 环的投射极限上局部自由”的形式群，建立了通常针对
域上形式群所陈述的一般理论。关于后者的更详细研究，请参阅 1964/65 年
Heidelberg--Strasbourg 代数几何讨论班。\footnote{编者只找到了 1965/66 年题为“线性代数群〔Linear Algebraic Groups〕”的这一讨论班，
其中并未出现形式群的概念；请改参阅 \textup{[De72]}。}

\section*{0. 伪紧环与伪紧模回顾}
\phantomsection
\label{sga3:VIIB:0}
\addcontentsline{toc}{section}{0. 伪紧环与伪紧模回顾}

本节包含若干技术性准备知识。我们在这里回顾并补充 \textup{[CA]}
（\emph{阿贝尔范畴〔Des cat\'egories ab\'eliennes〕}，Bull.\ Soc.\ Math.\ France
\textbf{90} (1962)）中的若干结果。

\medskip
\noindent\textbf{0.1.}
\phantomsection
\label{sga3:VIIB:0.1}
一个\emph{左伪紧环}是指带单位元的分离完备拓扑环，它具有由有限余长度的
左理想 \(\mathfrak l\) 构成的 \(0\) 的邻域基，即
\[
  \operatorname{length}_A(A/\mathfrak l)<+\infty.
\]
这里我们假定 \(A\) 交换，因此不必区分“左和右”。

特别地，商环 \(A/\mathfrak l\) 都是 Artin 环，而 \(A\) 等同于这些环
各自赋予离散拓扑后所得的拓扑投射极限。

完备 Noether 局部环 \((A,\mathfrak m)\) 显然是伪紧的。\footnote{这里赋予它 \(\mathfrak m\)-进拓扑。}

% 依据： Exp7B 原文件页码 2; 合并阅读版页码 515; 印刷页码 504.
\medskip
\noindent\textbf{0.1.1.}
\phantomsection
\label{sga3:VIIB:0.1.1}
\(A\) 的每个闭理想 \(I\) 都是包含它的所有开理想的交。\footnote{事实上，若 \(x\notin I\)，则存在开理想 \(\mathfrak l\) 使
\((x+\mathfrak l)\cap I=\varnothing\)；于是 \(I+\mathfrak l\) 是不包含
\(x\) 的开理想。在下文中，编者也明确说明了每个“闭极大”理想都是极大且开的。}
因此，每个闭极大理想都是开的。此外，若 \(\mathfrak l\) 是 \(A\) 的开理想，
则 \(A/\mathfrak l\) 的极大理想与 \(A\) 中包含 \(\mathfrak l\) 的极大理想
\(\mathfrak m\) 一一对应；后者因而既开又闭。因此，每个闭极大理想都是开极大理想
（从而是闭的），其逆命题显然成立。我们以 \(\Upsilon(A)\) 表示这些理想组成的集合。

若 \(\mathfrak l\) 是 \(A\) 的开理想，且
\(\mathfrak m\in\Upsilon(A)\)，则局部化
\((A/\mathfrak l)_{\mathfrak m}\) 在 \(\mathfrak m\) 包含
\(\mathfrak l\) 时是局部环，否则为零。由于 \(A/\mathfrak l\) 是 Artin 环，
它是有限多个局部环的直积；因此
\[
  A/\mathfrak l
  \simeq
  \prod_{\mathfrak m\in\Upsilon(A)}
    (A/\mathfrak l)_{\mathfrak m}.
\]
由此得到下列“典范”同构
\[
  A
  \simeq \varprojlim_{\mathfrak l}(A/\mathfrak l)
  \simeq \varprojlim_{\mathfrak l}
    \prod_{\mathfrak m}(A/\mathfrak l)_{\mathfrak m}
  \simeq \prod_{\mathfrak m}
    \varprojlim_{\mathfrak l}(A/\mathfrak l)_{\mathfrak m}
  \simeq \prod_{\mathfrak m}A_{\mathfrak m},
\]
其中
\[
  A_{\mathfrak m}
  =
  \varprojlim_{\mathfrak l}(A/\mathfrak l)_{\mathfrak m}.
\]
这个\emph{局部分量} \(A_{\mathfrak m}\) 是赋予离散拓扑的 Artin 局部环的
滤过投射极限。因此，它是一个局部环，并且对于投射极限拓扑是伪紧的。\footnote{注意，\(A_{\mathfrak m}\) 是 \(A\) 在极大理想 \(\mathfrak m\) 处的局部化。
事实上，\(A_{\mathfrak m}\) 的单位元 \(e_{\mathfrak m}\) 是 \(A\) 中不属于
\(\mathfrak m\) 的幂等元。由于 \(A=A_{\mathfrak m}\times B\)，其中
\(B=A(1-e_{\mathfrak m})\)，环 \(A_{\mathfrak m}\) 等同于局部化
\(A_{e_{\mathfrak m}}\)，从而也等同于 \(S^{-1}A\)，其中
\(S=A\setminus\mathfrak m\)。此外，\(e_{\mathfrak m}(1-e_{\mathfrak m})=0\)，
故 \(\mathfrak m\) 包含 \(1-e_{\mathfrak m}\)，从而也包含 \(B\)；
所以 \(\mathfrak m\) 等同于 \(\mathfrak mA_{\mathfrak m}\times B\)。}

\medskip
\noindent\textbf{0.1.2.}
\phantomsection
\label{sga3:VIIB:0.1.2}
设 \(r(A)\) 是 \(A\) 的所有开极大理想的交，即在下列等同下各理想
\(\mathfrak mA_{\mathfrak m}\) 的笛卡尔积
\[
  A=\prod_{\mathfrak m}A_{\mathfrak m}.
\]
对于 \(A\) 的每个开理想 \(\mathfrak l\)，\(r(A)\) 在
\(A/\mathfrak l\) 中的像包含于 \(A/\mathfrak l\) 的根中。
因此，这个像的某个幂为零，故对所有充分大的 \(n\) 都有
\(r(A)^n\subset\mathfrak l\)。于是序列 \(r(A)^n\) 趋于 \(0\)。

当 \(x\in r(A)\) 时，\(x^n\) 也如此。换言之，\(r(A)\) 的每个元素都是
拓扑幂零的，逆命题显然成立。因此，当 \(x\in r(A)\) 时，通项为
\[
  1+x+\cdots+x^n
\]
的序列收敛，并且收敛到 \(1/(1-x)\)。这表明 \(r(A)\) 是 \(A\) 的
Jacobson 根，即 \(A\) 的所有极大理想的交（参见\ Bourbaki，
\emph{代数〔Alg\`ebre〕}，第~8 章，\S6，定理~1）。\footnote{事实上，设 \(x\in r(A)\)。若 \(\mathfrak m\) 是不包含 \(x\) 的极大理想，
则存在 \(y\in A\) 使 \(1-xy\in\mathfrak m\)。但是 \(xy\in r(A)\)，
故 \(1-xy\) 可逆，矛盾。注意如下推论：若
\(\Upsilon(A)=\{\mathfrak m_1,\ldots,\mathfrak m_r\}\) 是有限集，
则这些 \(\mathfrak m_i\) 就是 \(A\) 的全部极大理想。}

\medskip
\noindent\textbf{注记。}\footnote{编者添加这些注记，是为了将 1.1 中形式谱 \(\operatorname{Spf}(A)\) 的定义
与 EGA~I，10.1.2 和 10.4.2 中的定义作比较。}

\noindent\textit{a)} 若 \(\mathfrak p\) 是 \(A\) 的开素理想，则因 \(A/\mathfrak p\) 是 Artin 环，
\(\mathfrak p\) 是极大理想。因此，\(\Upsilon(A)\) 是 \(A\) 的开素理想集。

\smallskip
\noindent\textit{b)} 每个 \(\mathfrak mA_{\mathfrak m}\) 都是 \(A_{\mathfrak m}\) 的定义理想，
即满足 \(I^n\) 趋于 \(0\) 的开理想 \(I\)
（参见\ EGA~0\textsubscript{I}，7.1.2）。因此，
\[
  \operatorname{Spec}
    (A_{\mathfrak m}/\mathfrak mA_{\mathfrak m}),
\]
配备拓扑环 \(A_{\mathfrak m}\) 后，是 EGA~I，10.1.2 意义下的仿射形式概形。

\smallskip
\noindent\textit{c)} 拓扑环 \(A\) 是 EGA~0\textsubscript{I}，7.1.2 意义下的容许环，
当且仅当 \(r(A)\) 是开的（从而是定义理想），而这又当且仅当
\(\Upsilon(A)\) 是有限集。这时，仿射形式概形
\[
  \operatorname{Spf}(A)
  =
  \operatorname{Spec}(A/r(A))
\]
（参见\ EGA~I，10.1.2）的底空间是赋予离散拓扑的 \(\Upsilon(A)\)；
其结构层在子集 \(E\subset\Upsilon(A)\) 上的截面构成环
\[
  \prod_{\mathfrak m\in E}A_{\mathfrak m}.
\]

\smallskip
\noindent\textit{d)} 设 \(A\) 是任意伪紧环。在下文 \SGARefLink{sga3:VIIB:1.1}{1.1} 中，
我们赋予空间 \(\Upsilon(A)\) 离散拓扑，以及在任意子集 \(E\) 上的截面环为
\[
  \prod_{\mathfrak m\in E}A_{\mathfrak m}.
\]
的环层。由 \textit{b)}，每个点都有一个为仿射形式概形的开邻域；
因而这定义了一个形式概形（EGA~I，10.4.2），记作 \(\operatorname{Spf}(A)\)。
这个形式概形若是仿射的，就必须拟紧，故 \(\Upsilon(A)\) 必须有限；
这时 \(\operatorname{Spf}(A)\) 与 EGA~I，10.1.2 中的定义一致。

% 依据： Exp7B 原文件页码 3; 合并阅读版页码 516; 印刷页码 505.
\medskip
\noindent\textbf{0.1.3.}
\phantomsection
\label{sga3:VIIB:0.1.3}
若 \(A\) 和 \(B\) 是伪紧环，则从 \(A\) 到 \(B\) 的\emph{同态}，
按定义是与加法、乘法及单位元相容的连续映射。这样的同态将 \(A\) 的拓扑幂零元素
映为 \(B\) 的拓扑幂零元素；因此，它将根 \(r(A)\) 映入根 \(r(B)\)。

% 下一断点： Exp7B 原文件页码 3, 合并阅读版页码 516,
% 再版印刷页码 505, 0.2 处。
